% GENERATED FILE. Edit canonical lessons, then run npm run generate:books.
% canonical-source-hash: fnv1a64:4fb5f0b7a6bf6f09
% canonical-lessons: JA-C358-yachin, JA-C358-ouya, JA-C358-kouza, JA-C358-furikomu, JA-C358-zandaka

\chapter{Renting a Room}
\label{ch:ja-renting-a-room}
\hlchaptermodality{358}

\begin{chapteropening}
\textbf{By the end of this chapter:} \emph{I can talk about rent and paying it.}

Complete the last lesson of chapter 358: I can talk about rent and paying it.
\end{chapteropening}

\section[\texorpdfstring{\ja{やちん} (yachin)}{やちん (yachin)}]{\ja{やちん} (yachin) — rent (for a home)}

\label{lesson:JA-C358-yachin}



\begin{quote}
Before the new one: say the Japanese for leeway, then the Japanese for to decide officially.
\end{quote}

\subsection*{You'll want to know: \ja{やちん}}
\textbf{\ja{やちん}} — \emph{yachin} — \textquotedblleft{}rent (for a home)\textquotedblright{}.

\begin{quote}
\textbf{\ja{やちんが} \ja{あがったら、ひっこします。}} — \emph{yachin ga agattara, hikkoshimasu} — \textquotedblleft{}If the rent goes up, I will move.\textquotedblright{}
\end{quote}

\begin{cousinweb}
\textbf{\ja{やちん}} (\emph{yachin}) is built from the native \emph{ya}, \textquotedblleft{}house\textquotedblright{}, a kun'yomi, and the Sino-Japanese \emph{chin}, \textquotedblleft{}fee\textquotedblright{}, an on'yomi borrowed from Chinese (Middle Chinese \emph{nrimH}) (KANJIDIC2; JMdict; Middle Chinese in Baxter's transcription).
\end{cousinweb}

\begin{grammarlens}[title={What you've built}]
One more word for saying what happens if.
\end{grammarlens}

\subsection*{Your turn}
\begin{itemize}
\raggedright
  \item \emph{Say it:} \emph{yachin}
  \item \emph{Say it:} \emph{yachin}, once more
  \item \emph{Say it:} \emph{yachin ga agattara, hikkoshimasu} — If the rent goes up, I will move.
  \item \emph{Recall:} say the Japanese for leeway, then the Japanese for to decide officially, then say \emph{yachin} again
\end{itemize}

\begin{tcolorbox}[breakable,colback=teal!4,colframe=teal!35!black,title={Before you move on}]
What does \ja{やちん} mean? (\textquotedblleft{}Rent (for a home)\textquotedblright{}.) Say it once more.
\end{tcolorbox}
\section[\texorpdfstring{\ja{おおや} (ōya)}{おおや (ōya)}]{\ja{おおや} (ōya) — a landlord; a landlady}

\label{lesson:JA-C358-ouya}



\begin{quote}
Before the new one: say the Japanese for to decide officially, then the Japanese for rent.
\end{quote}

\subsection*{You'll want to know: \ja{おおや}}
\textbf{\ja{おおや}} — \emph{ōya} — \textquotedblleft{}a landlord; a landlady\textquotedblright{}.

\begin{quote}
\textbf{\ja{へやが} \ja{こわれたら、おおやに} \ja{れんらくして} \ja{ください。}} — \emph{heya ga kowaretara, ōya ni renraku shite kudasai} — \textquotedblleft{}If something in the flat breaks, contact the landlord.\textquotedblright{}
\end{quote}

\begin{cousinweb}
\textbf{\ja{おおや}} (\emph{ōya}) is a native Japanese word (wago), inherited rather than borrowed from Chinese: it is built from \emph{ō}, \textquotedblleft{}big\textquotedblright{} and \emph{ya}, \textquotedblleft{}house\textquotedblright{}, and its kanji were attached later as kun'yomi (JMdict).
\end{cousinweb}

\begin{grammarlens}[title={What you've built}]
One more word for saying what happens if.
\end{grammarlens}

\subsection*{Your turn}
\begin{itemize}
\raggedright
  \item \emph{Say it:} \emph{ōya}
  \item \emph{Say it:} \emph{ōya}, once more
  \item \emph{Say it:} \emph{heya ga kowaretara, ōya ni renraku shite kudasai} — If something in the flat breaks, contact the landlord.
  \item \emph{Recall:} say the Japanese for to decide officially, then the Japanese for rent, then say \emph{ōya} again
\end{itemize}

\begin{tcolorbox}[breakable,colback=teal!4,colframe=teal!35!black,title={Before you move on}]
What does \ja{おおや} mean? (\textquotedblleft{}A landlord; a landlady\textquotedblright{}.) Say it once more.
\end{tcolorbox}
\section[\texorpdfstring{\ja{こうざ} (kōza)}{こうざ (kōza)}]{\ja{こうざ} (kōza) — a bank account}

\label{lesson:JA-C358-kouza}



\begin{quote}
Before the new one: say the Japanese for rent, then the Japanese for a landlord.
\end{quote}

\subsection*{You'll want to know: \ja{こうざ}}
\textbf{\ja{こうざ}} — \emph{kōza} — \textquotedblleft{}a bank account\textquotedblright{}.

\begin{quote}
\textbf{\ja{こうざが} \ja{あったら、おかねを} \ja{あずけて} \ja{ください。}} — \emph{kōza ga attara, okane o azukete kudasai} — \textquotedblleft{}If you have a bank account, please put the money in it.\textquotedblright{}
\end{quote}

\begin{cousinweb}
\textbf{\ja{こうざ}} (\emph{kōza}) is Sino-Japanese (kango): two Chinese characters read with their on'yomi, readings borrowed from Chinese. They are \emph{kō}, \textquotedblleft{}mouth\textquotedblright{} (Middle Chinese \emph{khuwX}) and \emph{za}, \textquotedblleft{}seat\textquotedblright{} (\emph{dzwaH}) (KANJIDIC2; Middle Chinese in Baxter's transcription).
\end{cousinweb}

\begin{grammarlens}[title={What you've built}]
One more word for saying what happens if.
\end{grammarlens}

\subsection*{Your turn}
\begin{itemize}
\raggedright
  \item \emph{Say it:} \emph{kōza}
  \item \emph{Say it:} \emph{kōza}, once more
  \item \emph{Say it:} \emph{kōza ga attara, okane o azukete kudasai} — If you have a bank account, please put the money in it.
  \item \emph{Recall:} say the Japanese for rent, then the Japanese for a landlord, then say \emph{kōza} again
\end{itemize}

\begin{tcolorbox}[breakable,colback=teal!4,colframe=teal!35!black,title={Before you move on}]
What does \ja{こうざ} mean? (\textquotedblleft{}A bank account\textquotedblright{}.) Say it once more.
\end{tcolorbox}
\section[\texorpdfstring{\ja{ふりこむ} (furikomu)}{ふりこむ (furikomu)}]{\ja{ふりこむ} (furikomu) — to pay into an account}

\label{lesson:JA-C358-furikomu}



\begin{quote}
Before the new one: say the Japanese for a landlord, then the Japanese for a bank account.
\end{quote}

\subsection*{You'll want to know: \ja{ふりこむ}}
\textbf{\ja{ふりこむ}} — \emph{furikomu} — \textquotedblleft{}to pay into an account\textquotedblright{}.

\begin{quote}
\textbf{\ja{らいげつに} \ja{なったら、やちんを} \ja{ふりこんで} \ja{ください。}} — \emph{raigetsu ni nattara, yachin o furikonde kudasai} — \textquotedblleft{}When next month comes, please pay the rent into the account.\textquotedblright{}
\end{quote}

Polite \textbf{\ja{ふりこみます}}, \emph{furikomimasu}; a godan verb. Past (polite): \textbf{\ja{ちちが} \ja{やちんを} \ja{ふりこみました。}}, \emph{chichi ga yachin o furikomimashita}.

\begin{cousinweb}
\textbf{\ja{ふりこむ}} (\emph{furikomu}) is a native Japanese word (wago), inherited rather than borrowed from Chinese: it is built from \emph{furi}, \textquotedblleft{}swinging (furu)\textquotedblright{} and \emph{komu}, \textquotedblleft{}to go into\textquotedblright{}, and its kanji were attached later as kun'yomi (JMdict).
\end{cousinweb}

\begin{grammarlens}[title={What you've built}]
One more word for saying what happens if.
\end{grammarlens}

\subsection*{Your turn}
\begin{itemize}
\raggedright
  \item \emph{Say it:} \emph{furikomu}
  \item \emph{Say it:} \emph{furikomu}, once more
  \item \emph{Say it:} \emph{raigetsu ni nattara, yachin o furikonde kudasai} — When next month comes, please pay the rent into the account.
  \item \emph{Recall:} say the Japanese for a landlord, then the Japanese for a bank account, then say \emph{furikomu} again
\end{itemize}

\begin{tcolorbox}[breakable,colback=teal!4,colframe=teal!35!black,title={Before you move on}]
What does \ja{ふりこむ} mean? (\textquotedblleft{}To pay into an account\textquotedblright{}.) Say it once more.
\end{tcolorbox}
\section[\texorpdfstring{\ja{ざんだか} (zandaka)}{ざんだか (zandaka)}]{\ja{ざんだか} (zandaka) — a bank balance}

\label{lesson:JA-C358-zandaka}



\begin{quote}
Before the new one: say the Japanese for a bank account, then the Japanese for to pay into an account.
\end{quote}

\subsection*{You'll want to know: \ja{ざんだか}}
\textbf{\ja{ざんだか}} — \emph{zandaka} — \textquotedblleft{}a bank balance\textquotedblright{}.

\begin{quote}
\textbf{\ja{ざんだかが} \ja{すくなく} \ja{なったら、おしえて} \ja{ください。}} — \emph{zandaka ga sukunaku nattara, oshiete kudasai} — \textquotedblleft{}If the balance gets low, please let me know.\textquotedblright{}
\end{quote}

\begin{cousinweb}
\textbf{\ja{ざんだか}} (\emph{zandaka}) is built from the Sino-Japanese \emph{zan}, \textquotedblleft{}remainder\textquotedblright{}, an on'yomi borrowed from Chinese (Middle Chinese \emph{dzan}), and the native \emph{taka}, \textquotedblleft{}amount\textquotedblright{}, voiced to \emph{daka} (KANJIDIC2; JMdict; Middle Chinese in Baxter's transcription).
\end{cousinweb}

\begin{grammarlens}[title={What you've built}]
One more word for saying what happens if.
\end{grammarlens}

\subsection*{Your turn}
\begin{itemize}
\raggedright
  \item \emph{Say it:} \emph{zandaka}
  \item \emph{Say it:} \emph{zandaka}, once more
  \item \emph{Say it:} \emph{zandaka ga sukunaku nattara, oshiete kudasai} — If the balance gets low, please let me know.
  \item \emph{Recall:} say the Japanese for a bank account, then the Japanese for to pay into an account, then say \emph{zandaka} again
\end{itemize}

\begin{tcolorbox}[breakable,colback=teal!4,colframe=teal!35!black,title={Before you move on}]
What does \ja{ざんだか} mean? (\textquotedblleft{}A bank balance\textquotedblright{}.) Say it once more.
\end{tcolorbox}
